\documentclass[sigconf]{acmart}

\settopmatter{printacmref=false}
\setcopyright{none}
\renewcommand\footnotetextcopyrightpermission[1]{}

\usepackage{amsmath}
\usepackage{booktabs}
\usepackage{float}

% --------------------------------------------------
% Paper information
% --------------------------------------------------

\title{NLP Assignment 1 Q2 Report}

\author{Muhammad Umer Siddiqui}
\email{ms08702@st.habib.pk}
\affiliation{%
  \institution{Dhanani School of Science and Engineering, Habib University}
  \city{Karachi}
  \country{Pakistan}
}

\author{Abdullah Shaikh}
\email{as09245@st.habib.pk}
\affiliation{%
  \institution{Dhanani School of Science and Engineering, Habib University}
  \city{Karachi}
  \country{Pakistan}
}

\author{Arqish Zaria}
\email{az09714@st.habib.pk}
\affiliation{%
  \institution{Dhanani School of Science and Engineering, Habib University}
  \city{Karachi}
  \country{Pakistan}
}

\begin{document}

\begin{abstract}
This report presents a preprocessing, word-frequency and part-of-speech
analysis of Jane Austen's \textit{Emma}, loaded from NLTK's Gutenberg
corpus. Preprocessing (em-dash spacing, lowercasing, an alphabetic filter
and stopword removal) reduced the text tokens and the vocabulary of unique words. The most frequent remaining words are dominated by character names and by modal verbs that NLTK's stopword list does not cover. POS tagging was run sentence by sentence on the unfiltered text, and the noun analysis
is dominated by character names and honorifics. The report also discusses how preprocessing choices affected the tagger.
\end{abstract}

\keywords{Natural Language Processing, Text Preprocessing, Word Frequency,
POS Tagging, NLTK}

\maketitle

% --------------------------------------------------
% Main content
% --------------------------------------------------

\section{Introduction}

The question required an analysis of an English text; preprocessing it and
comparing statistics before and after, finding the most frequent words,
running a part-of-speech (POS) tagger to analyse the nouns and discussing
how the preprocessing decisions affected the results. We used Jane
Austen's \textit{Emma}, loaded with NLTK's Gutenberg corpus reader
(\texttt{austen-emma.txt}).

The whole analysis is implemented in \texttt{q2\_final.py}. Besides printing the results, the script writes three output files: \texttt{top10\_words.txt}
(the ten most frequent words after preprocessing),
\texttt{top10\_nouns.txt} (the ten most frequent nouns) and
\texttt{unique\_nouns.txt} (every distinct noun found, sorted
alphabetically). All numbers in this report are taken from the output of
this script and from these files.

\section{Part A: Text Preprocessing}

The raw text was first tokenized as is with NLTK's \texttt{word\_tokenize},
with no modification, to obtain baseline statistics. The baseline total
token count includes every token produced by the tokenizer (punctuation
included). The baseline unique-word count only counts alphabetic tokens
(tokens for which \texttt{isalpha()} is true) and is case sensitive so
\textit{The} and \textit{the} count as two different words.

One normalization step was then applied before computing the
after preprocessing statistics every occurrence of \texttt{--} (Project
Gutenberg's rendering of an em dash, written with no surrounding
whitespace) was replaced with \texttt{ -- } surrounded by spaces. Without
this step, punctuation next to an unspaced em dash stays attached to the
neighbouring word. For example, in \textit{``...brought it.--Nobody...''}
the period stays fused to the word before the dash and the token becomes
\textit{``it.''} which would later fail the alphabetic filter and get mistagged. 

After this normalization, every token was lowercased and passed through two
filters: an alphabetic check (\texttt{isalpha()}) and a stopword filter
using NLTK's built-in English stopword list. The alphabetic check removes
punctuation, digits and any other non-word tokens. Note that it also
removes tokens containing an apostrophe or a period, such as contraction
fragments (\textit{n't}, \textit{'s}) and dotted abbreviations such as
\textit{Mr.}, so those tokens do not reach the frequency analysis later on.

\begin{table}[H]
\caption{Token and vocabulary counts before and after preprocessing}
\label{tab:partA}
\centering
\begin{tabular}{lrr}
\toprule
 & Total tokens & Unique words \\
\midrule
Before preprocessing & 191{,}855 & 7{,}378 \\
After preprocessing  & 70{,}171  & 6{,}831 \\
\bottomrule
\end{tabular}
\end{table}

Preprocessing removed 63.42\% of all tokens (191,855 to 70,171) but only
about 7.41\% of the unique words (7,378 to 6,831). The token count falls
sharply because punctuation and stopwords are the most repeated tokens in
the novel whereas each of them is only a single entry in the vocabulary.
The unique-word count falls by a much smaller amount because the words
removed are few in number and because lowercasing merges case variants
(such as \textit{The} and \textit{the}) into one entry which also reduces
the count slightly.

\section{Part B: Word Frequency Analysis}

This part identifies the ten most frequent words remaining after
preprocessing, computed with a frequency counter over the filtered token
list.

\begin{table}[H]
\caption{Top 10 most frequent words after preprocessing}
\label{tab:partB}
\centering
\begin{tabular}{clr}
\toprule
Rank & Word & Frequency \\
\midrule
1  & emma    & 865 \\
2  & could   & 837 \\
3  & would   & 819 \\
4  & miss    & 599 \\
5  & must    & 567 \\
6  & harriet & 506 \\
7  & much    & 486 \\
8  & said    & 484 \\
9  & one     & 447 \\
10 & weston  & 439 \\
\bottomrule
\end{tabular}
\end{table}

The number of words occurring exactly once (hapax legomena) is 2,682,
which is about 39.26\% of the 6,831 unique words left after preprocessing.

Two things stand out in this list. First, the character names \textit{emma},
\textit{harriet} and \textit{weston}, together with the address term
\textit{miss}, take four of the ten places, which is expected in a novel
built around a small set of characters and their social interactions.
Second, the modal verbs \textit{could}, \textit{would} and \textit{must}
survived stopword removal and occupy three of the top five places. NLTK's default English stopword list does not include these
modal auxiliaries so words that behave more like grammatical connectors
than content words still pass through the filter. The generic words
\textit{much}, \textit{said} and \textit{one} also appear in the list.

\section{Part C: Part-of-Speech Tagging}

Nouns were defined as any token tagged \texttt{NN}, \texttt{NNS},
\texttt{NNP} or \texttt{NNPS}, i.e. singular and plural common nouns
together with singular and plural proper nouns. Proper noun tags are
included because a large share of the nouns in this text are character
names and excluding them would have made the analysis mostly about generic
words rather than the people the story is about.

POS tagging was performed on the em-dash-normalized text \emph{before}
lowercasing or stopword removal rather than on the preprocessed token
list. The text was split into sentences with \texttt{sent\_tokenize} and
each sentence was tokenized and tagged separately with \texttt{pos\_tag}.
Tagging on the filtered, lowercased token list would remove exactly the two
signals the tagger relies on: capitalization (needed to recognize proper
nouns and titles) and sentence level context (needed to resolve ambiguous
words). The em-dash is also applied before tagging because a fused token such as \textit{``it.''} can be tagged as a noun even
though the underlying word is a pronoun.

In addition to the tag based filter, each noun token was required to match
the pattern of one or more alphabetic characters with an optional single
trailing period. This excludes Project Gutenberg formatting artifacts such
as underscore-delimited italics (\textit{``\_her\_''}) and stray bracket
tokens (\textit{``[''}, \textit{``]''}), while still allowing abbreviations
such as \textit{Mr.} and single-letter initials to pass through. All nouns
were lowercased after tagging so that capitalized and lowercase
occurrences of the same word (for example \textit{Time} and \textit{time})
are counted together.

\begin{table}[H]
\caption{Top 10 most frequent nouns (NN, NNS, NNP, NNPS)}
\label{tab:partC}
\centering
\begin{tabular}{clr}
\toprule
Rank & Noun & Frequency \\
\midrule
1  & mr        & 1{,}153 \\
2  & emma      & 857 \\
3  & mrs       & 699 \\
4  & miss      & 584 \\
5  & harriet   & 494 \\
6  & weston    & 439 \\
7  & thing     & 398 \\
8  & knightley & 389 \\
9  & elton     & 385 \\
10 & woodhouse & 311 \\
\bottomrule
\end{tabular}
\end{table}

The total number of noun occurrences is 31,317 and the number of distinct
nouns is 3,711.

The titles \textit{mr} and \textit{mrs} appear without their trailing period
because of an additional normalization step. The tokenizer sometimes splits
the period off an abbreviation into its own token (for example when the
abbreviation falls at a detected sentence boundary) so both the bare form
\textit{mr} and the dotted form \textit{mr.} can occur in the noun list.
For every noun token ending in a period, the code checks whether the same
word also occurs in the noun list without a period. If it does, the two
are merged into a single count under the bare spelling; if no bare form
exists (for example \textit{dr.} or \textit{sept.} in
\texttt{unique\_nouns.txt}), the period is kept. This is why the titles in
Table~\ref{tab:partC} report a single combined frequency (1,153 for
\textit{mr} and 699 for \textit{mrs}) rather than a count split across two
spellings of the same word.

Apart from these two titles, the list consists of six character names
(\textit{emma}, \textit{harriet}, \textit{weston}, \textit{knightley},
\textit{elton}, \textit{woodhouse}), the address term \textit{miss} and one
generic common noun, \textit{thing}.

\subsection{Residual Tagging Errors}
\label{sec:residual}

The full noun list in \texttt{unique\_nouns.txt} shows that the tagger is
not perfect. Besides genuine nouns, it contains some function words and
other non-nouns, for example \textit{a}, \textit{am}, \textit{are},
\textit{was}, \textit{were}, \textit{had}, \textit{her}, \textit{him},
\textit{my} and \textit{which}, as well as Roman numerals from chapter
headings (\textit{ii}, \textit{iii}, \textit{iv}, ... \textit{xix}) and
single-letter initials. These entries are part of the 3,711 distinct nouns
reported above, so that figure should be read as an upper bound on the
number of true nouns. None of them appears in the top ten of
Table~\ref{tab:partC}.

\section{Part D: Analysis}

\textbf{Effect of preprocessing on vocabulary and frequency.} Preprocessing
reduced the token count by about 63.42\% but the unique-word count by only
about 7.41\% (Table~\ref{tab:partA}). This gap is expected: a handful of
function words and punctuation marks such as \textit{the}, \textit{to},
\textit{and} and commas account for a large share of the raw tokens but
each is only one vocabulary entry so removing them shrinks the token count
far more than the number of word types.

\textbf{Effect of stopword removal on the most frequent words.} Stopword
removal substantially changed the frequency list: the function words that
would dominate an unfiltered count, such as \textit{the}, \textit{to} and
\textit{and} do not survive it. However, it is not a complete fix.
NLTK's default stopword list omits modal auxiliaries so \textit{could},
\textit{would} and \textit{must} still occupy three of the top five places
in Table~\ref{tab:partB}, and other generic words such as \textit{much},
\textit{said} and \textit{one} also remain. The most frequent word is
therefore not necessarily the most meaningful word and the choice of
stopword list has a direct effect on the frequency table. Also, because the
alphabetic filter removes dotted abbreviations, \textit{mr} and \textit{mrs}
are largely absent from the word list in Part B, although they are the most
frequent nouns in Part C.

\textbf{Types of nouns that dominate.} Excluding the two honorifics, the
list in Table~\ref{tab:partC} consists of the six character names, the
address term \textit{miss} and the generic noun \textit{thing}. This fits a
novel that is driven mostly by conversation and social relations among characters rather than by description of objects or
settings.

\textbf{Preprocessing decisions that affected POS tagging.} Tagging the
text before any lowercasing or stopword removal was a deliberate choice:
keeping capitalization let the tagger recognize character names and titles
such as \textit{Mr.} and \textit{Mrs.} as proper nouns and tagging
sentence by sentence preserved the sentence level context the tagger uses
to resolve ambiguous words. Lowercasing the nouns only after tagging let
case variants be merged into a single count without affecting the tagging
decision.

The clearest example of a preprocessing decision affecting tagging is the
em-dash handling. A fused token like \textit{``it.''}, produced when a
word and its period were followed directly by \texttt{--}, could be tagged
as a noun because of a Project Gutenberg formatting quirk and not because
of a grammatical decision by the tagger; the word \textit{it} on its own
is a pronoun. Adding spaces around the em dash before tokenizing removes
this source of error. The residual errors listed in Section~\ref{sec:residual} show that
some tagging mistakes still remain so a frequency list should always be
checked for unexpected entries after a tagging pass.

\section{Conclusion}

This report analysed \textit{Emma} with NLTK. Preprocessing reduced the
text from 191,855 to 70,171 tokens (63.42\%) and the vocabulary from 7,378
to 6,831 unique words (7.41\%). The most frequent remaining words are
character names and modal verbs that NLTK's stopword list does not remove,
and 2,682 words (39.26\% of the remaining vocabulary) occur only once. POS
tagging, performed sentence by sentence on the unfiltered text, found
31,317 noun occurrences across 3,711 distinct nouns, dominated by
honorifics and character names. The results show that preprocessing choices,
in particular the stopword list, capitalization handling and the treatment
of the em dash, have a direct effect on both the frequency statistics and
the quality of POS tagging.

\end{document}